% Stage 8 restored fragment — 2.3.2 Point electric dipole and the dyadic Green operator (archive lines 3953-4085)
\subsection{Point electric dipole and the dyadic Green operator}
\label{sec:ch23.signature-dipole}
Engineering reading (Hertzian chain). \emph{Physical setup:} a short wire or probe feed whose electrical size satisfies
\(k_0a\ll 1\) on a turntable range.
\emph{Objective:} verify rank-one far-zone export before multipole catalogs are applied
to larger apertures.
\emph{Expected intuition:} a transverse \(\sin\theta\) pattern with a single dominant
radiative channel.
\emph{Measurement:} watt closure on a sphere \(\Gamma_R\) and turntable invariance of the
total radiated power \cite{jackson1999}.
\emph{Misconception:} equating mesh unknowns in a numerical model with the radiative rank
of the physical source \cite{chu1948}.

On a turntable compact range, the dipole test is performed before any aperture catalog:
rotate the electrically small feed, integrate $|\E|^2$ on $\Gamma_R$, and confirm that the
total radiated power is invariant and that the far-zone pattern carries a single transverse
lobe \cite{hansen1988}. That watt closure is the laboratory observation that
precedes the dyadic formulas of Subsection~\ref{sec:ch231}; when it fails, the fault lies
in calibration, outgoing branch selection, or source admissibility---not in the multipole
catalog applied later \cite{stratton1941}.
The Hertzian dipole is the canonical image of the dyadic Green operator: a localized
source, a retarded kernel, a radiation map, and a far-zone angular spectrum linked in
one closed calculation \cite{stratton1941}. Short-wire or probe feeds at $k_0a\ll1$ on
a turntable should reproduce transverse $\sin\theta$ export and turntable-invariant watt
totals before multipole catalogs are applied to larger apertures
\cite{jackson1999}. Every compact feed should be checked against
this lowest-dimensional outgoing test at low $ka$ \cite{chu1948}.
The range observation fixed above---watt closure on $\Gamma_R$, turntable invariance,
and a single transverse lobe---is what a correct Green propagation must reproduce before
any multipole catalog is trusted on larger structures \cite{hansen1988}.
Only after that calibration does it pay to write the dyadic convolution in closed form:
replace the extended source $\Jimp(\mathbf{r}')$ by a rank-one impulse and read off the
field pair that every method-of-moments code must match at electrically small $k_0 a$
\cite{harrington2001}.
The Hertzian dipole is chosen as the canonical calibration source because it is the
smallest oriented radiator that still exercises the full dyadic structure: one moment
vector, one propagation direction, and a transverse far-zone pattern that every radiation
theory must reproduce \cite{jackson1999}. Compact feeds on turntable ranges,
probe references, and method-of-moments self-tests all reduce to this limit at
$k_0 a\ll 1$; when a code or measurement sequence cannot match the dipole field pair, no
larger aperture catalog built on the same Green operator should be trusted
\cite{hansen1988}.
The dipole moment \(\mathbf{p}=\int \Jimp\,d^{3}r'\) (units: A\,m, equivalently C\,m for
equivalent polarization sources) localizes the impressed current at the origin; $\mathbf{r}$
denotes the observation point and $\mathbf{r}'$ the source coordinate in the display below.
Substituting \(\Jimp(\mathbf{r}')=\mathbf{I}\,\delta(\mathbf{r}')\) into
Eq.~\eqref{eq:ch2.electric-dyadic-convolution} collapses the volume integral to a single
dyadic--vector product and yields the closed-form field pair
\begin{equation}
\label{eq:ch2.signature-dipole-field}
\E(\mathbf{r})
=
-\jj\omega\mu_0\,\overline{\mathbf{G}}_{E,0}(\mathbf{r},\mathbf{0})\cdot\mathbf{p},
\qquad
\H(\mathbf{r})=\frac{1}{\jj\omega\mu_0}\,\curl\E(\mathbf{r}),
\end{equation}
with outgoing dyad Eq.~\eqref{eq:ch2.free-space-dyadic-green}. This is
\(\mathcal{S}_\omega[\Jimp]\) in closed form; \(\mathcal{R}_\omega[\Jimp]\) is the
same pair after outgoing selection on the exterior domain
\cite{stratton1941}. Physically, Eq.~\eqref{eq:ch2.signature-dipole-field} states that
every Cartesian component of $\mathbf{p}$ excites a column of $\overline{\mathbf{G}}_{E,0}$,
so polarization coupling is already present in the rank-one source that calibrates compact
feeds \cite{stratton1941}.
\begin{figure}[t]
\centering
\ChOneFigBlock{%
\begin{tikzpicture}[ch1fig, node distance=5mm and 7mm]
 \node[ch1box] (p) {Dipole $\mathbf{p}$};
 \node[ch1box, right=7mm of p] (g) {$\overline{\mathbf{G}}_{E,0}$};
 \node[ch1box, right=7mm of g] (e) {$\E,\H$};
 \node[ch1box, right=7mm of e] (r) {$\mathcal{R}_\omega$};
 \node[ch1box, right=7mm of r] (a) {$a(\widehat{\mathbf{k}})$};
 \draw[ch1flow] (p) -- (g) -- (e) -- (r) -- (a);
\end{tikzpicture}%
}
\caption{Signature chain for a Hertzian dipole: impressed moment convolves with the
outgoing dyadic Green function, producing \((\E,\H)\); the radiation map extracts the
far-zone amplitude and angular spectrum.}
\label{fig:ch2.signature-hertzian-dipole}
\end{figure}
In the radiative far zone \(k_0 r\gg 1\), with \(r=|\mathbf{r}|\) the observation radius
in metres, the reactive \(R^{-2}\) and \(R^{-3}\) terms in
Eq.~\eqref{eq:ch2.free-space-dyadic-green} become negligible and
Eq.~\eqref{eq:ch2.signature-dipole-field} reduces to
Eq.~\eqref{eq:ch2.far-field-amplitude-integral} with
\begin{equation}
\label{eq:ch2.signature-dipole-far-field}
\E(\mathbf{r})\approx
-\jj\omega\mu_0\,\frac{e^{-\jj k_0 r}}{4\pi r}\,
\overline{\mathbf{I}}-\widehat{\mathbf{r}}\widehat{\mathbf{r}}\,\cdot\mathbf{p},
\end{equation}
so the transverse radiative image is fixed by \(\widehat{\mathbf{r}}\times(\widehat{\mathbf{r}}\times\mathbf{p})\).
This expression reveals that only transverse components survive at large distance. The far
field therefore forgets most local source details and retains only radiatively accessible
information---precisely the content a compact-range pattern measurement can resolve on
$\Gamma_R$ \cite{stratton1941}.
Power through sphere \(S_r\) follows Eq.~\eqref{eq:ch2.radiated-power-screen} with
\(P_{\mathrm{rad}}\propto k_0^{4}|\mathbf{p}|^{2}\) in the electrically small limit
\cite{stratton1941}. Experimentally, the angular spectrum is the decomposition of the
far-zone field into plane waves propagating in different directions $\widehat{\mathbf{k}}$:
each outgoing direction receives a distinct amplitude fixed by how the dipole moment
projects transversely onto that direction \cite{hansen1988}. Pattern measurements on a
spherical scanner therefore sample this angular distribution; synthesis codes invert the
same lateral Fourier image when reconstructing aperture fields from far-zone data
\cite{stratton1941}. The plane-wave angular spectrum of
Subsection~\ref{sec:ch291} reads
\begin{equation}
\label{eq:ch2.signature-dipole-angular-spectrum}
a(\widehat{\mathbf{k}})
=
-\jj k_0^{2}\,\sqrt{\mu_0/\varepsilon_0}\,
\big(\overline{\mathbf{I}}-\widehat{\mathbf{k}}\widehat{\mathbf{k}}\big)\cdot\mathbf{p},
\end{equation}
where \(a(\widehat{\mathbf{k}})\) is the far-zone plane-wave amplitude (V\,m\(^{-1}\))
associated with propagation direction \(\widehat{\mathbf{k}}\) (dimensionless unit vector)
and \(\widehat{\mathbf{k}}\) parametrizes the angular-spectrum parametrization of
Subsection~\ref{sec:ch291}. Equation~\eqref{eq:ch2.signature-dipole-angular-spectrum}
is the lateral Fourier image of the same export channel: each outgoing direction receives
transverse amplitude fixed by $(\overline{\mathbf{I}}-\widehat{\mathbf{k}}\widehat{\mathbf{k}})\cdot\mathbf{p}$.
One source, one kernel, one radiation map, one far field, one angular spectrum---every
abstract operator statement in Chapter~\ref{ch:02} reduces to this geometry when
\(\Jimp=\mathbf{p}\,\delta(\mathbf{r}')\) \cite{stratton1941}. The dipole test on
$\Gamma_R$ is therefore not a numerical afterthought; it is the lowest-dimensional
witness that $G$, $\overline{\mathbf{G}}_{E}$, $\mathcal{S}_\omega$, and
$\mathcal{R}_\omega$ describe the same hardware \cite{hansen1988}.
